\documentclass[sigconf,review]{acmart}

\usepackage{amsmath}
\let\Bbbk\relax
\usepackage{amssymb}
\usepackage{booktabs}
\usepackage{xcolor}
\usepackage{graphicx}
\usepackage{url}
\theoremstyle{acmdefinition}
\newtheorem{remark}{Remark}

\copyrightyear{2027}
\acmYear{2027}
\acmDOI{10.1145/0000000.0000000}

\begin{document}

\title{Formal Verification of Lightning Network Multi-Hop HTLC Routing: From Concrete Models to Arbitrary Path Lengths}

\author{First Author}
\affiliation{
  \institution{University XYZ}
  \city{City}
  \country{Country}
}
\email{email@university.xyz}

\author{Second Author}
\affiliation{
  \institution{University XYZ}
  \city{City}
  \country{Country}
}
\email{email@university.xyz}

\author{Third Author}
\affiliation{
  \institution{Institute of Technology, ABC}
  \city{City}
  \country{Country}
}
\email{email@abc.edu}

\begin{abstract}
The Lightning Network (LN) relies on Hashed Time-Locked Contracts (HTLCs) and
complex channel state machines to enable scalable, multi-hop blockchain
payments. Prior symbolic analyses have targeted either two-party HTLC settings
or generic message-forwarding abstractions; the interaction between the channel
revocation lifecycle, multi-hop routing, and block-height timing has not been
mechanized as a single symbolic model. We present such a model in the Tamarin
prover, structured as a modular package of five theories totalling \textbf{81
machine-checked lemmas}, covering the two-party channel lifecycle, multi-hop
HTLC routing with fees, staggered-CLTV timing, and an $N$-hop generalisation.

Our central finding is a consequence of that integration. We identify a
\emph{linear single-spend discipline}: two independent soundness gaps---HTLC
output duplication and revoked-commitment replay---share one root cause,
namely modeling one-shot on-chain resources as Tamarin \emph{persistent} facts
rather than \emph{linear} ones, and both are repaired by the same
linear-token remedy. The second instance is visible only when revocation and
routing are modeled together, which narrower prior models structurally could
not observe; we propose the resulting invariant as a reusable modeling
discipline for symbolic payment-channel analyses. We further prove
economic-soundness theorems (end-to-end conservation, strictly positive fees)
and expose a griefing DoS; establish \emph{assumption minimality} empirically,
showing that removing the intermediary-claim restriction yields a concrete
59-step counterexample; and reproduce the known wormhole (agent-skipping)
attack as a fidelity check on the model. Finally, we identify Tamarin's
origin-analysis recursion on self-referential signed messages as a barrier to
unbounded verification, and give a proof-layering technique---machine-checked
local authentication plus a pen-and-paper refinement theorem---that transfers
core routing invariants to paths of arbitrary length.
\end{abstract}

\begin{CCSXML}
<ccs2012>
<concept>
<concept_id>10002978.10003014.10003015</concept_id>
<concept_desc>Security and privacy~Formal methods and theory of security</concept_desc>
<concept_significance>500</concept_significance>
</concept>
<concept>
<concept_id>10002978.10003029.10003030</concept_id>
<concept_desc>Security and privacy~Cryptography</concept_desc>
<concept_significance>300</concept_significance>
</concept>
</ccs2012>
\end{CCSXML}

\ccsdesc[500]{Security and privacy~Formal methods and theory of security}
\ccsdesc[300]{Security and privacy~Cryptography}

\keywords{Formal verification, Tamarin prover, Lightning Network, Payment Channels, HTLC, Blockchain, Layer-2 scaling}

\maketitle

% ==============================================================
\section{Introduction}
% ==============================================================

Blockchain technologies have revolutionized decentralized finance, but their inherent reliance on global consensus creates a fundamental scalability bottleneck. To achieve the transaction throughput required for global adoption, Layer-2 (off-chain) scaling solutions have become essential. The Lightning Network (LN)~\cite{poon2016bitcoin} is the most prominent of these, enabling fast, low-fee payments by moving the bulk of transactions off-chain and settling only disputes on the underlying blockchain~\cite{gudgeon2020sok}.

At the heart of the Lightning Network are Payment Channels and Hashed Time-Locked Contracts (HTLCs)~\cite{bolt2024}. HTLCs enable multi-hop payments across a network of untrusted intermediaries by conditioning the release of funds on the knowledge of a secret preimage, combined with staggered absolute timelocks (CLTV). Alongside this, the channel lifecycle relies on a complex revocation and punishment mechanism to prevent participants from publishing outdated, fraudulent channel states~\cite{maffei2024lecture}.

While the conceptual design of LN multi-hop routing is elegant, the interaction between cryptographic primitives, block-height timing constraints, and off-chain state machines is highly non-trivial. Subtle edge cases—such as early-timeout races where an incoming channel refunds before an outgoing channel is redeemed—can compromise the safety of honest intermediaries. Furthermore, security in these systems is heavily dependent on liveness assumptions: honest nodes must monitor the network and act before their HTLCs expire.

Existing formal treatments have addressed parts of this stack in isolation. Symbolic analyses in Tamarin have targeted two-party HTLC constructions such as cross-chain swaps~\cite{boyd2020blockchain} and contingent payments~\cite{bursuc2019contingent}; generic symbolic frameworks have studied path integrity and agent-skipping in message-forwarding protocols~\cite{smith2022agentskipping}; and model checking with TLA+ has verified balance security for Lightning channels and bounded-length multi-hop payments~\cite{grundmann2025modelchecking}. What has not been mechanized is the \emph{combination}: the two-party revocation and punishment lifecycle together with multi-hop HTLC routing, per-hop fees, and staggered CLTV timing, in one symbolic model under a full Dolev--Yao adversary. To our knowledge, the model we present is the first Tamarin formalization to integrate these layers.

That integration is not merely a matter of coverage; it is what makes our central finding observable. We identify a category error---modeling one-shot on-chain resources as \emph{persistent} rather than \emph{linear} Tamarin facts---that manifests twice, once in the HTLC layer and once in the revocation layer. The second instance involves a rule gated by revocation state being fired repeatedly along a routing execution, and is therefore invisible to models that contain only one of the two layers.

To conduct this analysis we use the Tamarin prover~\cite{meier2013tamarin} under a full Dolev--Yao adversary~\cite{dolev1983security}. We explicitly model integrity and safety; we do not model payment privacy or unlinkability (a confidentiality property requiring onion routing and observational-equivalence proofs), stated here as an explicit exclusion. Comprehensively verifying the Lightning Network introduces severe methodological and tooling hurdles. First, the protocol inherently requires the combination of Dolev-Yao cryptographic reasoning, natural-number induction (for CLTV arithmetic), and unbounded block clocks. As we demonstrate, attempting to verify these aspects in a single monolithic theory causes state-of-the-art provers to exhaust available memory (OOM). Second, while prior analyses often fix the routing path to a specific length, a natural question arises: do these safety properties hold for paths of arbitrary length ($N$ hops)?

In this paper, we present a comprehensive, modular formal verification of Lightning Network multi-hop HTLC routing. Our contributions are as follows:
\begin{enumerate}
    \item A \textbf{linear single-spend discipline}: a category error in which one-shot on-chain outputs are modeled as persistent facts, identified in two independent instances (HTLC output duplication and revoked-commitment replay) and repaired uniformly by structural linear-resource guards, replacing global axioms. We propose the resulting invariant---\emph{every one-shot on-chain output is a linear fact}---as a reusable modeling discipline for symbolic payment-channel analyses.
    \item A modular, fully machine-checked model integrating the two-party channel lifecycle (including revocation and punishment) with multi-hop HTLC routing and fees, verified across \textbf{81 lemmas} in five theories. This integration is the precondition for the finding above.
    \item An empirical assumption-minimality result proving that removing the intermediary-claim restriction yields a concrete 59-step counterexample, together with a two-sided certificate that the corresponding liveness assumption is load-bearing.
    \item Economic-soundness theorems (end-to-end value conservation, strictly positive per-hop fees) and a reachable griefing DoS, plus a reproduction of the known wormhole (agent-skipping) attack~\cite{malavolta2019anonymous, smith2022agentskipping} as a fidelity check on the model.
    \item A map of the prover's tractability boundaries, identifying the precise bounding factors (the signing/nat/clock OOM triple, signed-message $N$-hop routing, and reusable channels), and a proof-layering technique---machine-checked local authentication plus a pen-and-paper refinement theorem---for $N$-hop invariants.
\end{enumerate}

% ==============================================================
\section{Related Work}
% ==============================================================

\subsection{Formal Verification of Payment Channels}
Early formalizations with idealized channel closures~\cite{decker2015fast, malavolta2017concurrency, miller2019sprites} successfully proved payment atomicity and balance conservation but omitted the revocation and punishment mechanisms required to secure real-world implementations.

Prior symbolic analyses have applied Tamarin to HTLC-based cross-chain swaps~\cite{boyd2020blockchain} and zero-knowledge contingent payments over a public ledger~\cite{bursuc2019contingent}. These models target two-party settings and lack multi-hop routing, channel lifecycle revocation, or staggered CLTV arithmetic. Our work explicitly bridges this gap within a unified, automated framework targeting the multi-hop setting.

\subsection{Agent-Skipping and Path Integrity}
Smith et al.~\cite{smith2022agentskipping} develop a general symbolic framework, in the multiset-rewriting setting, for \emph{agent-skipping} attacks on message-forwarding protocols, introducing path-integrity goals and analyzing several protocols including payment channel networks. The wormhole attack on HTLC routing~\cite{malavolta2019anonymous} is an instance of agent-skipping, and their framework establishes it as a symbolic reachability result. Our treatment of the wormhole is therefore \emph{not} a novel attack finding but a fidelity check: we confirm that our integrated model reproduces the known attack, and additionally tie it to the value layer by quantifying the theft to exactly the bypassed routing fees within a model that also carries channel revocation and settlement. The complementary distinction is one of scope: Smith et al.\ abstract the forwarding protocol to study path integrity generically, whereas we model the concrete Lightning stack---revocation, punishment, fees, CLTV---in order to reason about properties that depend on those layers.

\subsection{Model Checking and Computational Treatments}
Grundmann and Hartenstein~\cite{grundmann2025modelchecking} formalize Lightning in TLA+ and model-check balance security, using stepwise refinement and a time abstraction to make single channels and multi-hop payments checkable separately, scaling to payments over up to four hops with two concurrent payments. Their approach is complementary to ours along two axes. Methodologically, model checking explores a bounded state space with an explicit adversary encoding, whereas symbolic verification in Tamarin reasons over unbounded sessions under an equational Dolev--Yao model; consequently their analysis handles concurrency that we do not, while ours handles unbounded message-derivation capabilities and, via the abstraction of Section~\ref{sec:nhop}, unbounded path length. Substantively, their refinement establishes that the current Lightning specification satisfies balance security; our contribution is orthogonal, concerning a modeling pitfall in the symbolic encoding of one-shot on-chain outputs.

Kiayias and Litos~\cite{kiayias2020composable} provide a composable, game-based (UC) security treatment of the LN at the two-party channel layer; it does not machine-check multi-hop HTLC routing. Our symbolic approach achieves a high degree of automation---proving 81 lemmas---while explicitly covering multi-hop routing under a full Dolev--Yao adversary. Prior work frequently treats the liveness requirement of intermediaries as an implicit assumption. The Tamarin analysis of HTLC swaps by Boyd et al.~\cite{boyd2020blockchain} surfaced a timing-related hidden assumption (stable relative chain growth), which parallels our finding that liveness assumptions are strictly load-bearing; we provide a mechanized, two-sided certificate for this assumption in a model that includes the full channel lifecycle.

\subsection{Tamarin in Blockchain Protocols}
The Tamarin prover~\cite{meier2013tamarin} is a standard tool for verifying cryptographic protocols, successfully applied to industrial standards such as TLS 1.3~\cite{cremers2017comprehensive}. In the blockchain domain, Tamarin has been applied to HTLC swaps and contingent payments~\cite{boyd2020blockchain, bursuc2019contingent}. Existing Tamarin applications typically verify fixed-round protocols. Applying Tamarin to stateful, long-running protocols with unbounded loops introduces unique challenges. We do not merely restrict the model to avoid OOM; we rigorously diagnose the root causes, pinpointing Tamarin's origin-analysis recursion on self-referential signed terms as the fundamental barrier to verifying generic $N$-hop routing.

% ==============================================================
\section{Protocol Overview and Threat Model}
% ==============================================================

\subsection{Payment Channel Lifecycle}
A payment channel is established via an on-chain funding transaction. Once opened, parties transact off-chain by exchanging signed commitment transactions. To prevent a malicious party from broadcasting an outdated state, LN employs a revocation mechanism. When transitioning to a new state, parties exchange revocation secrets. If a revoked state is published, the other party can unilaterally punish the cheater and claim the entirety of the channel funds, protected by a relative timelock (CSV). Channels can be closed cooperatively or unilaterally.

\subsection{Multi-Hop HTLC Routing}
To pay a recipient without a direct channel, LN uses HTLCs to route payments across intermediate nodes. Assume a path $Sender \rightarrow F_1 \rightarrow \cdots \rightarrow F_n \rightarrow Receiver$.
\begin{enumerate}
    \item \textbf{The Lock Round (Setup):} The Receiver generates a random secret preimage $x$ and sends its hash $y = H(x)$ to the Sender. The Sender adds an HTLC to their channel with $F_1$, locking funds conditional on providing $x$ before an absolute deadline. This cascades down, with each forwarder taking a routing fee.
    \item \textbf{The Release Round (Settlement):} The Receiver reveals $x$ to claim the final HTLC. $F_n$ uses $x$ to claim the HTLC from $F_{n-1}$. This cascades backward to the Sender.
\end{enumerate}

\subsection{Threat Model and Security Goals}
We adopt the standard Dolev-Yao adversary model~\cite{dolev1983security}, representing a network attacker with full control over the communication medium but who cannot break cryptographic primitives. We explicitly model long-term key compromise as a distinct adversarial capability.

\textbf{The Liveness Assumptions.} Security degrades into a liveness requirement. Our model encodes four timing restrictions:
\begin{itemize}
\item \textbf{T1} (\texttt{RedeemBeforeTimeout});
\item \textbf{T2a} (\texttt{ForwardTimeGap});
\item \textbf{T2b} (\texttt{HonestPartiesActBefore\-IncomingTimeout});
\item \textbf{T3} (\texttt{IntermediaryMustClaim}).
\end{itemize}
As we demonstrate in Section~\ref{sec:verification}, T3 is strictly load-bearing. This mirrors the real-world necessity of LN watchtowers~\cite{mccorry2019pisa}. Our core security goals translate directly to verified lemmas: Payment Atomicity, Intermediary Safety, Routing Causality \& Authentication, Timing Safety, and Value Conservation.

% ==============================================================
\section{Formal Preliminaries}
\label{sec:prelim}
% ==============================================================

We fix notation for Tamarin's multiset-rewriting semantics~\cite{meier2013tamarin}, which
the remainder of the paper uses to state both the modeling discipline of
Section~\ref{sec:discipline} and the refinement theorem of Section~\ref{sec:nhop}.

\paragraph{Terms and facts.}
Let $\mathcal{T}$ be the set of terms over a signature with an equational theory
$E$ (here: hashing, signing, and natural numbers). A \emph{fact} is
$F(t_1,\dots,t_k)$ with $t_i \in \mathcal{T}$. The set of fact symbols is
partitioned as $\Sigma_{\mathrm{fact}} = \Sigma_{\mathrm{lin}} \uplus
\Sigma_{\mathrm{per}}$ into \emph{linear} and \emph{persistent} symbols; we write
persistent facts with a leading bang, e.g.\ $!\mathsf{RevokedSecret}(\cdot)$.
A \emph{state} $S$ is a finite multiset of ground facts.

\paragraph{Rules and transitions.}
A rule is $ri = l \overset{a}{\longrightarrow} r$ with $l,a,r$ multisets of facts
($a$ the \emph{action} facts). Writing $\mathrm{lin}(l)$ and $\mathrm{per}(l)$ for
the linear and persistent premises, the transition relation is
\[
S \xrightarrow{\;a\;} \bigl(S \setminus^{\#} \mathrm{lin}(l)\bigr) \cup^{\#} r
\qquad\text{provided}\quad
\mathrm{lin}(l) \subseteq^{\#} S \ \wedge\ \mathrm{per}(l) \subseteq S ,
\]
where $\setminus^{\#},\cup^{\#},\subseteq^{\#}$ are multiset operations. The
asymmetry is the point: linear premises are \emph{removed} from the state, while
persistent premises are merely \emph{checked} and remain available. A
\emph{trace} is the sequence $\mathit{tr} = [a_1,\dots,a_n]$ of action multisets
of an execution from the empty state; $\mathrm{traces}(R)$ denotes the traces of
a rule set $R$. Properties are first-order formulas over
$\mathit{F}(\bar t)@i$ with timepoint ordering $i \lessdot j$, interpreted over
traces in the standard way; $R \models_{\forall} \varphi$ means all traces
satisfy $\varphi$ and $R \models_{\exists} \varphi$ that some trace does.

\paragraph{Restrictions.}
A restriction $\rho$ filters the trace set: verification is carried out over
$\mathrm{traces}_{\rho}(R) = \{\mathit{tr} \in \mathrm{traces}(R) \mid \mathit{tr} \models \rho\}$.
Restrictions are therefore \emph{assumptions}, not theorems, which is why
Section~\ref{sec:verification} tests whether one of them is load-bearing.

% ==============================================================
\section{Formal Modeling Methodology}
% ==============================================================

\subsection{The Curse of Monolithism and Modular Decomposition}
A naive approach is to construct a single, monolithic theory containing the channel lifecycle, HTLC routing logic, CLTV block-height arithmetic, and value conservation. This is fundamentally intractable. The combination of the signature equational theory, natural-number induction, and an unbounded block clock causes Tamarin's search space to explode. Crucially, this was tested, not assumed—any two of the three ingredients remain tractable; only the full triple OOM-kills the prover.

The model is therefore split so that no single theory carries all three. The architecture is summarized in Table~\ref{tab:theories}.

\begin{table*}[t]
\centering
\small
\caption{The five-theory modular architecture (81 lemmas total).}
\label{tab:theories}
\begin{tabular}{@{}llcl@{}}
\toprule
\textbf{Theory File} & \textbf{Aspect Modeled} & \textbf{Lemmas} & \textbf{Builtins} \\
\midrule
\texttt{multihop.spthy} & Lifecycle, 3-hop routing, value/fees, wormhole & 43 & Hashing, Signing, Nat \\
\texttt{multihop\_nhop.spthy} & N-hop generalisation (generic linear-fact forward) & 25 (+4 comm.) & Hashing, Signing, Nat \\
\texttt{Clock.spthy} & Block-clock lifecycle \& staggered-CLTV claim windows & 9 & Natural-numbers \\
\texttt{cltv\_blocks.spthy} & Pure CLTV block arithmetic (delta positivity) & 3 & Natural-numbers \\
\texttt{timeout.spthy} & Early-timeout race (T2b removed) & 1 & --- \\
\bottomrule
\end{tabular}
\end{table*}

\subsection{Tactic Engineering and Reproduction}
All results were produced under a strict test-first discipline: a property is claimed only after Tamarin reports it. Per-lemma invocation is used on the large theory to avoid OOM. The flag \texttt{-{}-stop-on-trace=seqdfs} rescues the heavy existence witnesses but hangs the natural-number arithmetic theories, so it is applied selectively.

% ==============================================================
\section{Verification of the Core Model}
\label{sec:verification}
% ==============================================================

Across the five theories detailed in Table~\ref{tab:theories}, we successfully prove 81 lemmas. Table~\ref{tab:lemmas} maps the lemma groups in our primary concrete theory.

\begin{table*}[t]
\centering
\small
\caption{Lemma groups in \texttt{multihop.spthy} (43 total).}
\label{tab:lemmas}
\begin{tabular}{@{}lll@{}}
\toprule
\textbf{Group} & \textbf{Class} & \textbf{Representative Lemmas} \\
\midrule
Invoice / preimage auth & Safety & \texttt{Invoice\_Authenticates\_Settlement} \\
Lock causality & Safety & \texttt{Forward1\_Requires\_Offer} \\
Release causality & Safety & \texttt{Fulfill\_Requires\_Forward2} \\
Atomicity / non-vacuity & Safety/Reach. & \texttt{Payment\_Atomicity\_Under\_Liveness} \\
Value \& fees & Safety & \texttt{EndToEnd\_Value\_Conservation} \\
Refund / griefing & Reachable & \texttt{Honest\_Intermediary\_Refunds} \\
Wormhole & Attack & \texttt{Wormhole\_Steals\_Exactly\_The\_Fees} \\
Intermediary safety & Safety & \texttt{Intermediary\_Never\_Loses\_Under\_Liveness} \\
Channel lifecycle & Safety & \texttt{No\_Punishment\_Without\_Cheating} \\
\bottomrule
\end{tabular}
\end{table*}

\subsection{The Linear Single-Spend Discipline (Axiom to Structure)}
\label{sec:discipline}

We now state formally the pitfall informally described above, using the notation
of Section~\ref{sec:prelim}.

\begin{definition}[One-shot resource]
\label{def:oneshot}
A fact symbol $\mathsf{Res}$ \emph{denotes a one-shot resource} if, in the
intended protocol semantics, for every ground instantiation $\bar t$ the
privilege represented by $\mathsf{Res}(\bar t)$ may be exercised at most once.
On-chain outputs---a funding output, an HTLC output, the output of a revoked
commitment---are one-shot resources.
\end{definition}

\begin{proposition}[Persistent gating admits unbounded exercise]
\label{prop:bug}
Let $ri = \bigl(P \uplus \{!\mathsf{Res}(\bar t)\}\bigr) \overset{a}{\longrightarrow} r$
be a rule in which the only premise representing the resource is the persistent
fact $!\mathsf{Res}(\bar t)$, and suppose the remaining premises $P$ are
persistent or otherwise re-derivable. Then for every $n \in \mathbb{N}$ there is
a trace of $R$ containing at least $n$ occurrences of the action $a$ under the
same instantiation $\bar t$.
\end{proposition}

\begin{proof}
By the transition rule, firing $ri$ removes only $\mathrm{lin}(l)$ from the
state. Since $!\mathsf{Res}(\bar t) \in \Sigma_{\mathrm{per}}$ it is checked but
not removed, and by hypothesis $P$ is re-derivable; hence the enabling condition
$\mathrm{lin}(l) \subseteq^{\#} S \wedge \mathrm{per}(l) \subseteq S$ holds again
in the successor state. Firing $ri$ $n$ times yields the required trace.
\end{proof}

Proposition~\ref{prop:bug} is exactly the defect we found twice. In Instance~1
the resource is the channel output hosting an HTLC and $a$ is an HTLC-add; in
Instance~2 the resource is a revoked commitment, $!\mathsf{Res} =
{!\mathsf{RevokedSecret}}$, and $a = \mathsf{PublishRevoked}$, so the protocol
admits arbitrarily many punishments for a single cheat.

\begin{definition}[Linear single-spend token]
\label{def:token}
A \emph{single-spend token} for a one-shot resource is a linear fact symbol
$\mathsf{Tok} \in \Sigma_{\mathrm{lin}}$ together with a discipline requiring
(i) every rule exercising the resource to take $\mathsf{Tok}(\bar t)$ as a linear
premise, and (ii) $\mathsf{Tok}(\bar t)$ to be produced at most once per $\bar t$
across all rules.
\end{definition}

\begin{proposition}[Structural single spend]
\label{prop:fix}
Under Definition~\ref{def:token}, for every trace and every $\bar t$ the action
$a$ occurs at most once with instantiation $\bar t$.
\end{proposition}

\begin{proof}
Let $p$ be the number of productions and $c$ the number of consumptions of
$\mathsf{Tok}(\bar t)$ in a trace. As $\mathsf{Tok}$ is linear, each firing of a
rule exercising the resource removes one copy from the state, so $c \le p$;
by (ii) $p \le 1$, hence $c \le 1$. By (i) each occurrence of $a$ under $\bar t$
is such a firing, so $a$ occurs at most once.
\end{proof}

Concretely, $\mathsf{Free}(\mathit{ptr})$ is produced once by
$\mathsf{Lock\_Funds\_And\_Open}$ (one per channel open) and consumed by every
HTLC-add, and
$\mathsf{RevokedCommitmentUnspent}(\mathit{ptr},A,B,n,\mathit{commit})$ is
produced once by $\mathsf{Revoke\_Old\_Secret}$ and consumed by
$\mathsf{Publish\_Revoked}$:
\[
\footnotesize
\begin{array}{l}
\mathsf{Publish\_Revoked_A}:\\[2pt]
\bigl[\ !\mathsf{ChannelRecord}(\mathit{ptr},A,B),\ \ !\mathsf{ValidOnChain}(\mathit{ptr}),\\
\ \ \ !\mathsf{RevokedSecret}(A,n,s,c),\\
\ \ \ \mathsf{RevokedCommitmentUnspent}(\mathit{ptr},A,B,n,c)\ \bigr]\\[2pt]
\quad\overset{\ \mathsf{PublishRevoked}(A,n,c),\ \mathsf{Cheat}(A,n)\ }{\longrightarrow}\\[2pt]
\bigl[\ \mathsf{OnChainCheat}(A,B,n,c),\ \ \mathsf{RevealedSecret}(A,n,s,c)\ \bigr]
\end{array}
\]
Persistent \emph{knowledge} of the revocation secret is retained via
$!\mathsf{RevokedSecret}$; only the one-shot \emph{spend} is linear. This
separation is the content of the discipline.

\begin{remark}[Why re-verification is required]
\label{rem:mono}
Adding a premise shrinks the enabling condition, so
$\mathrm{traces}(R') \subseteq \mathrm{traces}(R)$. Hence every
$\forall$-property is preserved: $R \models_{\forall} \varphi \Rightarrow R'
\models_{\forall} \varphi$. The converse fails for $\exists$-properties, since a
witness may be removed. We therefore re-ran the full theory rather than
appealing to monotonicity; all 43 lemmas re-verify (262\,s, Tamarin 1.8).
\end{remark}

\medskip
\noindent\textbf{Instance 1: HTLC output duplication.} The forwarding rules were gated only by a public HTLC-offer message. A Dolev--Yao replay of that offer could re-fire an honest forward, locking a second HTLC on the same channel output with an adversary-chosen amount. We replaced a global \texttt{OneHTLCPerPtr} axiom with the structural guard of Definition~\ref{def:token}: each channel open emits one linear $\mathsf{Free}(\mathit{ptr})$ token, consumed by every HTLC-add.

\medskip
\noindent\textbf{Instance 2: revoked commitment replay.} The punishment rules were gated only by the persistent $!\mathsf{RevokedSecret}$ fact, and are thus an instance of Proposition~\ref{prop:bug}: the punishment rule could fire arbitrarily many times for the same revoked state.

We stress that Instance~2 is a consequence of integration rather than of scale.
It resides at the junction of the revocation lifecycle and the routing
execution: the offending rule is gated by revocation state yet re-fired along a
routing trace. A model containing only the two-party channel layer, or only an
abstracted forwarding layer, does not exhibit the trace that exposes it.

\textbf{The reusable invariant.} Definitions~\ref{def:oneshot}--\ref{def:token}
and Propositions~\ref{prop:bug}--\ref{prop:fix} yield a mechanical audit
question for symbolic payment-channel models: \emph{for each fact denoting a
one-shot on-chain output, is there a linear token consumed when it is
exercised?} Persistent facts may encode durable knowledge (keys, revealed
secrets, channel identity), but never a spend opportunity.

\subsection{Wormhole Attack: A Fidelity Check}
The wormhole attack~\cite{malavolta2019anonymous} occurs when two colluding path members relay the preimage out-of-band, skipping an honest intermediary who forwards value but earns no fee. It is an instance of the agent-skipping class formalized symbolically by Smith et al.~\cite{smith2022agentskipping}, and we claim no novelty for exhibiting it. We include it as a fidelity check: a model of multi-hop routing that did \emph{not} admit the wormhole would be suspect. In the symbolic model, the Dolev--Yao network is the colluders' side channel. We machine-check it twice: \texttt{Wormhole\_Fee\_Theft\_Reachable} (51 steps) exhibits a trace where the payment settles yet the middle hop is never redeemed and no intermediary earns a fee, with no key compromise; and \texttt{Wormhole\_Steals\_Exactly\_The\_Fees} (81 steps) quantifies the theft to precisely the path fees \texttt{\%f1 + \%f2}, tying the attack to the value layer of a model that simultaneously carries channel revocation and settlement. Known mitigations (AMHL, PTLCs) are out of scope.

\subsection{Economic Soundness and Griefing}
Three all-traces theorems: \texttt{EndToEnd\_Value\_Conservation} (the sender locks the receiver's amount plus a strictly positive total fee; no value created); \texttt{Fee\_Strictly\_Positive\_Hop1/2} (every hop's fee is $> 0$—structural, because the natural-number sort has no \texttt{\%0}); and, as a reachability result, \texttt{Honest\_Intermediary\_Refunds} (an honest intermediary induced to lock capital then refund on timeout—the standard PCN griefing DoS).

\subsection{Timing Safety and Assumption Minimality}
\texttt{cltv\_blocks} derives the CLTV-delta invariant from concrete block numbers; \texttt{Clock} adds block-clock lifecycle safety; and \texttt{timeout} exhibits the early-timeout race that reappears once T2b is removed.

A natural question is whether the intermediary-claim restriction T3 is redundant given T1 and T2b. We tested it by removing T3 and re-running the suite: \texttt{Intermediary\_Never\_Loses\_Under\_Liveness} is falsified, with Tamarin finding a concrete 59-step counterexample. T2b orders the two timeouts but does not force the node to sweep the incoming HTLC once its outgoing one is redeemed; that action is exactly what T3 encodes. We have empirically demonstrated that T3 is \textbf{load-bearing} for intermediary safety. We tested the removal of T3 only, and accordingly make no claim regarding the individual minimality of T1, T2a, or T2b.

% ==============================================================
\section{\texorpdfstring{Generalizing to Arbitrary Path Length ($N$-Hops)}{Generalizing to Arbitrary Path Length (N-Hops)}}
\label{sec:nhop}
% ==============================================================

Replacing the hardcoded forward rules with one generic rule was attempted in several encodings. A generic signed-message forward causes Tamarin to OOM even for a 1-hop witness due to self-referential origin-analysis recursion.

To bypass this, we abstract the network layer into an internal Tamarin linear fact—\texttt{Route(prev, me, ptr, y)}—handed directly from one hop to the next. Because the HTLC is no longer a self-referential signed term, Tamarin's origin analysis no longer recurses. This yields \texttt{multihop\_nhop.spthy} (25 active lemmas; 4 \texttt{[use\_induction]} lemmas are commented out as deliberate tractability boundaries because backward chaining through unbounded \texttt{Route} facts causes divergence).

Fees are enforced structurally in the N-hop model as a firing guard (\texttt{Eq(\%vIn, \%vOut + \%fee)}), meaning end-to-end conservation over $N$ hops is a sum of $N$ fees—not finitely statable in one Tamarin lemma. We note also that the abstraction internalises the network, and therefore deliberately does \emph{not} exhibit the wormhole: the side channel the attack exploits is absent by construction, which is why that result is established in the concrete model.

\subsection{Refinement Theorem}

Internalising the network removes Dolev--Yao power, so the abstract model $A$
has strictly fewer behaviours than the concrete model $C$; abstract safety
therefore does not by itself imply concrete safety. We make the transfer precise.

\begin{definition}[Shared alphabet and projection]
\label{def:proj}
Let
\[
\small
\Sigma = \{\,\mathsf{ForwardHTLC},\ \mathsf{Redeemed},\
\mathsf{Refunded},\ \mathsf{LtkCompromised}\,\}
\]
be the safety-relevant action symbols, emitted with identical arities by both
theories. The projection
$\pi : \mathrm{traces}(C) \to (2^{\Sigma})^{*}$ retains every $\Sigma$-action
verbatim, deletes all network and cryptographic events
($\mathsf{Out}, \mathsf{In}, \mathsf{KU}$, signature checks), and relabels the
position-indexed $\mathsf{HTLCForwarded}_i$ to the generic
$\mathsf{HTLCForwarded}$.
\end{definition}

\begin{definition}[\textsc{Auth}]
\label{def:auth}
$C$ satisfies \textsc{Auth} if every forwarding step is preceded by a genuine
offer or forward from the \emph{adjacent} previous hop, unless that hop's key is
compromised; for the first position,
\[
\small
\begin{array}{l}
\forall F_1,F_2,S,p,y,v_{\mathrm{in}},v_{\mathrm{out}},i.\
 \mathsf{HTLCForwarded}_1(F_1,F_2,S,p,y,v_{\mathrm{in}},v_{\mathrm{out}})@i \Rightarrow\\
\quad (\exists p',j.\ \mathsf{HTLCOffered}(S,F_1,p',y,v_{\mathrm{in}})@j \wedge j \lessdot i)\\
\quad \vee\ (\exists k.\ \mathsf{LtkCompromised}(S)@k \wedge k \lessdot i),
\end{array}
\]
and symmetrically for the second position.
\end{definition}

\textsc{Auth} is machine-checked in $C$ as \texttt{Forward1\_Requires\_Offer}
(266 steps) and \texttt{Forward2\_Requires\_Forward1} (274 steps). Crucially it
quantifies over a \emph{single adjacent pair} and never over the path or $N$;
this locality is what makes it provable, and its failure for a generic signed
rule is the boundary reported in Section~\ref{sec:nhop}.

\begin{lemma}[Per-hop simulation]
\label{lem:sim}
Assume $C \models_{\forall} \textsc{Auth}$. If $S \xrightarrow{a} S'$ is a
forwarding step of $C$, then $\pi$ of that step is a legal step of $A$.
\end{lemma}

\begin{proof}[Proof sketch]
An honestly authenticated incoming HTLC is precisely what licenses the abstract
$\mathsf{Route}(\mathit{prev},\mathit{me},\mathit{ptr},y)$ hand-off, so the
premises of the abstract $\mathsf{Forward\_HTLC}$ are met. The remaining case, a
forged or replayed forward, is excluded by \textsc{Auth} except under key
compromise, which both models expose identically via $\mathsf{LtkCompromised}
\in \Sigma$ and which $\pi$ preserves.
\end{proof}

\begin{theorem}[Soundness of the abstraction]
\label{thm:refine}
Let $\varphi$ be a safety property over $\Sigma$ with $A \models_{\forall}
\varphi$ for all path lengths $N$, and assume $C \models_{\forall} \textsc{Auth}$.
Then $C \models_{\forall} \varphi$.
\end{theorem}

\begin{proof}[Proof sketch]
By induction on the number of forwarding steps of $\mathit{tr} \in
\mathrm{traces}(C)$. The base case is immediate. For the step, Lemma~\ref{lem:sim}
extends $\pi(\mathit{tr})$ by one legal abstract step, so $\pi(\mathit{tr}) \in
\mathrm{traces}(A)$. Since $\pi$ preserves $\Sigma$-actions verbatim and
$\varphi$ is over $\Sigma$, any violation of $\varphi$ in $\mathit{tr}$ appears
in $\pi(\mathit{tr})$, contradicting $A \models_{\forall} \varphi$.
\end{proof}

\begin{remark}[Scope of mechanization]
\label{rem:scope}
Both hypotheses of Theorem~\ref{thm:refine} are machine-checked---\textsc{Auth}
in $C$, and $\varphi$ in $A$ for unbounded $N$---but the composition is
\textbf{pen-and-paper}: Tamarin can neither quantify over $N$ nor relate two
theories, for the same reason that key-management meta-theorems are argued
externally. Soundness rests on the locality of \textsc{Auth}
(Definition~\ref{def:auth}), which we confirmed by inspecting every forward
rule: each references only its adjacent pair. This is the honest boundary of our
verification.
\end{remark}

% ==============================================================
\section{Discussion: What Failed and Tractability Boundaries}
% ==============================================================

We record negative results explicitly—they are a core methodological contribution.
\begin{itemize}
    \item \textbf{The Three-Ingredient OOM.} Signing + natural-number induction + an unbounded block clock together exhaust the prover; any two are fine. This is the reason for the modular split.
    \item \textbf{Signed-Message $N$-Hop Routing.} Origin analysis on the self-referential \texttt{'htlc'} term does not terminate, even for reachability, even with a fresh ID.
    \item \textbf{A Rejected Restriction.} An initial fix \texttt{OneRedeemPerPtr} (``at most one redeem per output'') was tested and rejected: it falsified the honest witness \texttt{Multihop\_Payment\_Possible}, because the two \texttt{Redeemed} events in an honest run are the two endpoints' views of one hop, not two spends. The correct guard was \texttt{OneHTLCPerPtr} (later made structural).
    \item \textbf{A Caught False Positive.} An exists-trace probe sharing only the payment hash appeared to show ``value creation''. Tightening it to a single linked payment chain falsified it—the apparent attack was invoice-hash reuse across two unrelated offers, not a soundness bug.
    \item \textbf{Reusable Channels (Boundary).} Stage 1 (structural single-HTLC-per-output via the \texttt{Free(ptr)} slot) ships and re-verifies 43/43. Stage 2 (return \texttt{Free(ptr)} on resolution) is a clean tractability boundary: all-traces safety lemmas verify ($\sim$45 s each), but every exists-trace witness fails to converge. Stage 1 is therefore the shipped model.
\end{itemize}

% ==============================================================
\section{Conclusion and Future Work}
% ==============================================================

We delivered a modular, machine-checked model that integrates the Lightning two-party channel lifecycle with multi-hop HTLC routing and fees (81 lemmas over five theories). That integration surfaced a \emph{linear single-spend discipline}: two independent soundness gaps arising from encoding one-shot on-chain outputs as persistent facts, repaired uniformly by linear tokens, and generalizable into an auditable modeling invariant. We further established the minimality of the T3 timing assumption empirically, proved economic-soundness theorems, exposed a griefing DoS, reproduced the known wormhole attack as a fidelity check, and mapped where the prover stops scaling.

The protocol-level safety results confirm expectations—consistent with the composable cryptographic treatment of Lightning~\cite{kiayias2020composable} and the model-checking results of Grundmann and Hartenstein~\cite{grundmann2025modelchecking}, which our symbolic model complements at the routing layer. The durable contribution here is methodological—a modeling discipline that catches a class of silent encoding errors, together with an account of what it takes to keep Tamarin tractable on a stateful, time-locked, arithmetic protocol, in a regime where every claim (positive or negative) was executed rather than asserted.

Future work includes testing whether the persistent-versus-linear pitfall is latent in other published symbolic models of payment channels and blockchain protocols, which would establish the discipline as a class-level finding rather than a single-model observation; extending the framework to Point Time-Locked Contracts (PTLCs); and bridging to computational verification. Our value conservation proofs verify strict amount invariants, but do not enforce game-theoretic economic rationality (e.g., defending against timelock bribery), which remains outside symbolic reasoning.

% ==============================================================
\begin{thebibliography}{9}
\bibitem{bolt2024}
Lightning Network community.
2024.
BOLT: Basis of Lightning Technology (Lightning Network Specifications).
\url{https://github.com/lightning/bolts}.

\bibitem{boyd2020blockchain}
Colin Boyd, Kristian Gjøsteen, and Shuang Wu.
2020.
A Blockchain Model in Tamarin and Formal Analysis of Hash Time Lock Contract.
In \emph{2nd Workshop on Formal Methods for Blockchains (FMBC)}, 5:1--5:13.

\bibitem{bursuc2019contingent}
Sergiu Bursuc and Steve Kremer.
2019.
Contingent Payments on a Public Ledger: Models and Reductions for Automated Verification.
In \emph{European Symposium on Research in Computer Security (ESORICS)}. Springer.

\bibitem{cremers2017comprehensive}
Cas Cremers, Marko Horvat, Jonathan Hoyland, Sam Scott, and Thyla van der Merwe.
2017.
A Comprehensive Symbolic Analysis of TLS 1.3.
In \emph{ACM SIGSAC Conference on Computer and Communications Security (CCS)}. ACM.

\bibitem{decker2015fast}
Christian Decker and Roger Wattenhofer.
2015.
A Fast and Scalable Payment Network with Bitcoin Duplex Micropayment Channels.
In \emph{Stabilization, Safety, and Security of Distributed Systems (SSS)}. Springer, 3--18.

\bibitem{dolev1983security}
Danny Dolev and Andrew C. Yao.
1983.
On the Security of Public Key Protocols.
\emph{IEEE Transactions on Information Theory} 29, 2 (1983), 198--208.

\bibitem{grundmann2025modelchecking}
Matthias Grundmann and Hannes Hartenstein.
2025.
Model Checking the Security of the Lightning Network.
arXiv:2505.15568. (Supersedes arXiv:2307.02342.)

\bibitem{gudgeon2020sok}
Lewis Gudgeon, Pedro Moreno-Sanchez, Stefanie Roos, Patrick McCorry, and Arthur Gervais.
2020.
SoK: Layer-Two Blockchain Protocols.
In \emph{Financial Cryptography and Data Security (FC)}. Springer, 201--226.

\bibitem{kiayias2020composable}
Aggelos Kiayias and Orfeas Stefanos Thyfronitis Litos.
2020.
A Composable Security Treatment of the Lightning Network.
In \emph{IEEE 33rd Computer Security Foundations Symposium (CSF)}. IEEE, 334--349.

\bibitem{maffei2024lecture}
Matteo Maffei.
2024.
\emph{Cryptocurrencies, Lecture 11: Scalability (Payment Channels and Payment-Channel Networks)}. Lecture Notes.

\bibitem{malavolta2017concurrency}
Giulio Malavolta, Pedro Moreno-Sanchez, Aniket Kate, Matteo Maffei, and Srivatsan Ravi.
2017.
Concurrency and Privacy with Payment-Channel Networks.
In \emph{ACM SIGSAC Conference on Computer and Communications Security (CCS)}. ACM, 455--471.

\bibitem{malavolta2019anonymous}
Giulio Malavolta, Pedro Moreno-Sanchez, Clara Schneidewind, Aniket Kate, and Matteo Maffei.
2019.
Anonymous Multi-Hop Locks for Blockchain Scalability and Interoperability.
In \emph{Network and Distributed System Security Symposium (NDSS)}. The Internet Society.

\bibitem{mccorry2019pisa}
Patrick McCorry, Surya Bakshi, Iddo Bentov, Sarah Meiklejohn, and Andrew Miller.
2019.
Pisa: Arbitration Outsourcing for State Channels.
In \emph{ACM Conference on Advances in Financial Technologies (AFT)}. ACM, 16--30.

\bibitem{meier2013tamarin}
Simon Meier, Benedikt Schmidt, Cas Cremers, and David Basin.
2013.
The TAMARIN Prover for the Symbolic Analysis of Security Protocols.
In \emph{Computer Aided Verification (CAV)}. Springer, 696--701.

\bibitem{miller2019sprites}
Andrew Miller, Iddo Bentov, Surya Bakshi, Ranjit Kumaresan, and Patrick McCorry.
2019.
Sprites and State Channels: Payment Networks that Go Faster Than Lightning.
In \emph{Financial Cryptography and Data Security (FC)}. Springer.

\bibitem{poon2016bitcoin}
Joseph Poon and Thaddeus Dryja.
2016.
The Bitcoin Lightning Network: Scalable Off-Chain Instant Payments.
\url{https://lightning.network/lightning-network-paper.pdf}. Whitepaper.

\bibitem{smith2022agentskipping}
Zach Smith, Hugo Jonker, Sjouke Mauw, and Hyunwoo Lee.
2022.
Modelling Agent-Skipping Attacks in Message Forwarding Protocols.
arXiv:2201.08686.
\end{thebibliography}

\end{document}
